% How to reproduce the measurements and extend them to more configurations (input by
% MDWF_MULTIGRID_NOTE.tex with \resultstrue).

% ============================================================
\section{Reproducing and extending the measurements}
\label{sec:reproduce}
\begingroup\sloppy

All numbers of Sec.~\ref{sec:results} come from one configuration per volume; the
most useful extensions are more configurations, a larger volume and dynamical
ensembles. This section gives what is needed to run them.

\paragraph{Code.} RutticeQCD, \url{https://github.com/atsanroger/RutticeQCD}, branch
\texttt{cuda-bridge} (the commit of each result is given with it; later commits give
the same iterations, up to the two-iteration differences between builds noted in
Sec.~\ref{sec:massscan}). It needs Rust (stable) and, for the GPU, CUDA~12:
{\small
\begin{verbatim}
  export CUDA_HOME=/usr/local/cuda
  cargo build --release --features cuda-bridge
  ./target/release/ruttice mg_24.yaml
\end{verbatim}}
The GPU backend runs on one GPU. The parameter files are in the style of Bridge++
and are documented in \texttt{docs/USAGE.md}; gauge configurations are read in the
Bridge++ text and binary formats, NERSC, ILDG and SciDAC.

\paragraph{The mass scan (Table~\ref{tab:massscan}, Fig.~\ref{fig:massscan}).} The
parameter files are \texttt{docs/mg\_background/runs/mg\_24.yaml} (the multigrid) and
\texttt{runs/cgne\_24.yaml} (the reference, double or mixed precision); the comments
in them say how to switch to setup~B and to a point source. Each run prints a line
{\small
\begin{verbatim}
  seed 91: GCR 26 iterations on the device; 2875 coarse applications;
      worst |Dx-b|^2 2e-18; 24.9 s; ... device 19.91 GB
\end{verbatim}}
and, for the multigrid, \texttt{setup on the device: 393 s}. These, per mass and
setup, are what Table~\ref{tab:massscan} lists. On an A30 one configuration takes about
$1.5$ hours for the multigrid (seven masses, both setups; the setup dominates) and
about $40$ minutes for the two CGNE scans.

\paragraph{Memory and volume.} On $24^3\times48$ the multigrid needs $19.9$\,GB of
device memory (CGNE $6.9$\,GB in double, $10.1$\,GB mixed), roughly proportional to the
volume: $32^3\times64$ needs about $63$\,GB (an $80$\,GB GPU), $48^3\times96$ more
than one GPU has. The aggregates, $4^4$, must divide the lattice, and each level
below the coarse lattice must divide the extent of the level above
(\texttt{MultiGrid\_Level2}, \texttt{\_Level3}; on $32^3\times64$, for instance,
super-blocks $2$ and $2$).

\paragraph{The spectral quantities (Secs.~\ref{sec:fov}--\ref{sec:pseudo}).} The
example \texttt{preconditioned\_fov\_device} computes, with exact coarse solves, the
field of values, the parts of the cycle, the eigenvalues with their condition
numbers, and $\sigma_{\min}(z-\tilde A\mathcal M)$ on contours, e.g.
{\small
\begin{verbatim}
  cargo run --release --features cuda-bridge \
      --example preconditioned_fov_device -- \
      16 16 binary conf.bin 1.0 0.001 1 --ratio 10 \
      --arnoldi --far 1,0 --pseudo --contour 1,0,0.6,16
\end{verbatim}}
(arguments: $L$, $T$, format, file, $M_5$, $m$, stout $0/1$). It fits $16^4$ in $24$\,GB
and takes about an hour there; $24^3\times48$ does not fit. The CPU examples
\texttt{preconditioned\_fov} and \texttt{coarse\_hermiticity} need no GPU.

\paragraph{What would help most.}
\begin{enumerate}
\item The mass scan of Table~\ref{tab:massscan} on three to five further
$24^3\times48$ configurations, for the spread of the iterations and of the
speed-up.
\item The same on a larger volume ($32^3\times64$ on an $80$\,GB GPU), for the volume
dependence of the multigrid iterations and of the setup.
\item The mass scan on a dynamical ($2+1$-flavor) ensemble, e.g.\ that of
Ref.~\cite{KCM2024}, where the spectra of Sec.~\ref{sec:spectra} were taken.
\item The spectral run above on further $16^4$ configurations, for how the outliers
of setup~S and the pseudospectra vary between configurations.
\end{enumerate}
For each run, the result lines above and the commit of the build are enough; the
logs contain everything else.
\endgroup
